% ============================================================
%  CHAPITRE 3 — CONCEPTION ET ARCHITECTURE DE LA SOLUTION
%  Version détaillée, alignée sur la structure réelle du projet BFPME
%  (FastAPI · Option C RandomForest/XGBoost · SHAP · LLM Ollama ·
%   Monte Carlo · IFRS 9)
%
%  Ce fichier est COMPILABLE seul (preamble minimal identique au
%  rapport principal). Pour l'intégrer au rapport, supprimer le
%  preamble + \begin{document}/\end{document} et conserver à partir
%  de \chapter{Conception et architecture de la solution}.
%
%  NB : les diagrammes UML existants ne sont PAS modifiés ; ils sont
%  simplement référencés via \includegraphics (placeholders).
% ============================================================
\documentclass[a4paper,12pt,oneside]{report}

\usepackage[utf8]{inputenc}
\usepackage[T1]{fontenc}
\usepackage[french]{babel}
\usepackage[top=2.5cm, bottom=2.5cm, left=3cm, right=2.5cm]{geometry}
\usepackage{setspace}\onehalfspacing
\usepackage{lmodern}\usepackage{microtype}
\usepackage[dvipsnames,svgnames,table]{xcolor}
\definecolor{EspritRed}{RGB}{180, 0, 0}
\definecolor{DarkGray}{RGB}{50, 50, 50}
\usepackage{graphicx}\usepackage{tikz}\usepackage{pgffor}
\usepackage{titlesec}
\titleformat{\chapter}[display]
  {\normalfont\centering}
  {\rule{\linewidth}{2pt}\\[6pt]{\large\bfseries CHAPITRE \thechapter}\\\rule{\linewidth}{1pt}}
  {12pt}{\Huge\bfseries\MakeUppercase}[\vspace{4pt}\rule{\linewidth}{0pt}]
\titlespacing*{\chapter}{0pt}{-30pt}{40pt}
\titleformat{\section}{\normalfont\Large\bfseries}{\thesection}{1em}{}
\titleformat{\subsection}{\normalfont\large\bfseries}{\thesubsection}{1em}{}
\titleformat{\subsubsection}{\normalfont\normalsize\bfseries}{\thesubsubsection}{1em}{}
\titlespacing*{\section}{0pt}{16pt}{8pt}
\titlespacing*{\subsection}{0pt}{14pt}{6pt}
\titlespacing*{\subsubsection}{0pt}{10pt}{4pt}
\usepackage{fancyhdr}\pagestyle{fancy}\fancyhf{}
\fancyhead[L]{\small\textit{\leftmark}}\fancyhead[R]{\small PFE 2024--2025}
\fancyfoot[C]{\thepage}\renewcommand{\headrulewidth}{0.4pt}
\usepackage[colorlinks=true,linkcolor=black,citecolor=black,urlcolor=black]{hyperref}
\usepackage{booktabs}\usepackage{tabularx}\usepackage{longtable}
\usepackage{multirow}\usepackage{array}\usepackage{enumitem}
\setlist[itemize,1]{label=\textbullet, leftmargin=2em}
\setlist[itemize,2]{label=---, leftmargin=3em}
\usepackage{listings}
\lstset{
  basicstyle=\ttfamily\footnotesize, breaklines=true, frame=single,
  backgroundcolor=\color{gray!10},
  numbers=left, numberstyle=\tiny\color{gray}, numbersep=5pt,
  keywordstyle=\bfseries\color{EspritRed}, commentstyle=\itshape\color{gray},
  captionpos=b
}
\usepackage{amsmath, amssymb}
\usepackage{caption}\usepackage{subcaption}
\captionsetup{labelsep=endash, font=small, labelfont=bf, justification=centering}

\newcolumntype{Y}{>{\raggedright\arraybackslash}X}

\begin{document}
\setcounter{chapter}{2} % => ce chapitre sera numéroté « CHAPITRE 3 »

% ════════════════════════════════════════════════════════════
\chapter{Conception et architecture de la solution}

\section*{Introduction}

Après avoir posé, dans les chapitres précédents, le cadre métier du financement
des PME et le socle scientifique du projet, ce troisième chapitre est consacré à
la \textbf{conception et à l'architecture technique} de la plateforme effectivement
réalisée. Il fait le pont entre les besoins exprimés et l'implémentation : il décrit
l'architecture logique multi-couches, la structure concrète du projet (organisation
du code), les choix technologiques retenus, la conception des données, l'architecture
de l'API et des services métier, l'orchestration de l'agent IA, le module
d'explicabilité, le moteur de simulation Monte Carlo, la modélisation UML, ainsi
que la planification des sprints et des releases.

L'objectif de ce chapitre est double : justifier les décisions d'architecture au
regard des exigences fonctionnelles et non fonctionnelles (performance, modularité,
auditabilité, explicabilité), et fournir une vue suffisamment précise du système
pour en garantir la reproductibilité et la maintenabilité.

% ────────────────────────────────────────────────────────────
\section{Vision architecturale globale}

\subsection{Présentation générale de la plateforme}

La plateforme est conçue selon une \textbf{architecture modulaire en couches}, qui
sépare strictement les responsabilités : ingestion et validation des données,
prédiction de la probabilité de défaut (PD), explicabilité (SHAP), simulation de
scénarios et de stress (Monte~Carlo), orchestration décisionnelle et interaction en
langage naturel. Cette séparation des préoccupations
(\textit{separation of concerns}) répond directement aux exigences non
fonctionnelles de maintenabilité et d'évolutivité formulées au chapitre~1 : chaque
composant peut être testé, remplacé ou retraîné indépendamment des autres.

Le système est structuré autour de trois grands ensembles :

\begin{itemize}
  \item un \textbf{backend applicatif} exposant une API REST (FastAPI), qui
        orchestre la chaîne validation $\rightarrow$ prédiction $\rightarrow$
        explication $\rightarrow$ scénario $\rightarrow$ synthèse en langage naturel ;
  \item un \textbf{noyau analytique} (artefact de modèle sérialisé, moteur SHAP,
        moteur Monte~Carlo, ingénierie de variables) découplé de la couche HTTP ;
  \item une \textbf{interface analyste} (tableau de bord web statique) consommant
        l'API et restituant prédictions, explications et simulations sous forme
        graphique.
\end{itemize}

\subsection{Architecture logique multi-couches}

L'architecture s'organise en \textbf{cinq couches} complémentaires, chacune assurant
une fonction précise du pipeline décisionnel. La figure~\ref{fig:archi_couches}
en présente la vue d'ensemble.

\begin{figure}[h!]
  \centering
  % Remplacer par votre schéma d'architecture existant (non modifié)
  \includegraphics[width=0.95\textwidth]{architecture_multicouches}
  \caption{Architecture logique multi-couches de la plateforme}
  \label{fig:archi_couches}
\end{figure}

\subsubsection{Couche collecte et ingestion des données}

Cette couche centralise l'acquisition, la \textbf{validation} et le formatage des
données du dossier de crédit PME en provenance des systèmes de la banque. Le service
de validation (\texttt{validation\_service}) vérifie la présence des variables
attendues par le modèle, signale les champs manquants (qui seront imputés) et les
champs non utilisés, sans jamais bloquer la chaîne : la robustesse prime, les
manques étant gérés en aval par imputation.

\subsubsection{Couche modélisation et prédiction}

Elle héberge l'artefact entraîné (préprocesseur + modèle) chargé et mis en cache au
démarrage, et expose l'estimation de la PD pour chaque dossier soumis. C'est le
c\oe ur prédictif : application de l'ingénierie de variables, transformation par le
préprocesseur, puis \texttt{predict\_proba} et classement du risque
(\textit{low / medium / high}).

\subsubsection{Couche explicabilité et traçabilité}

Elle calcule les contributions \textbf{SHAP} (\texttt{TreeExplainer}) sur la matrice
de variables transformées et restitue, pour chaque prédiction, les facteurs qui
augmentent ou diminuent le risque. Chaque explication est structurée (valeur de base,
contributions ordonnées, direction) pour être auditable.

\subsubsection{Couche agent IA et orchestration}

Elle coordonne l'enchaînement des traitements (validation, prédiction, explication,
scénario), applique les seuils de confiance et gère la composition des résultats en
une réponse unique. C'est le niveau qui matérialise la logique d'\textit{agent} :
décider quels outils invoquer en fonction de la requête de l'analyste.

\subsubsection{Couche interaction LLM et interface analyste}

Elle expose une interface conversationnelle permettant aux analystes d'interroger le
système en langage naturel. Un grand modèle de langage (LLM) reçoit \textbf{uniquement}
les sorties calculées par les couches précédentes et produit une synthèse narrative,
sans jamais recalculer ni inventer de valeur numérique.

\subsection{Flux de données et interactions entre composants}

Les composants communiquent via des interfaces standardisées (modèles Pydantic
typés). Le flux nominal, depuis la soumission d'un dossier jusqu'à la restitution
d'une explication narrative, garantit la cohérence et la traçabilité de bout en bout.
La figure~\ref{fig:flux_donnees} en donne la représentation.

\begin{figure}[h!]
  \centering
  \includegraphics[width=0.9\textwidth]{flux_donnees}
  \caption{Flux de données de bout en bout du pipeline décisionnel}
  \label{fig:flux_donnees}
\end{figure}

% ────────────────────────────────────────────────────────────
\section{Structure du projet}

L'organisation du dépôt traduit directement l'architecture en couches : le backend
applicatif est isolé du noyau de données, des notebooks de modélisation et du module
réglementaire IFRS~9. Le listing~\ref{lst:arbo} présente l'arborescence simplifiée
du projet.

\begin{lstlisting}[caption={Arborescence simplifiée du projet}, label={lst:arbo}]
BFPME/
+-- backend/
|   +-- app/
|   |   +-- main.py            # App FastAPI, montage des routers + frontend
|   |   +-- config.py          # Configuration Pydantic (variables .env)
|   |   +-- schemas.py         # Modeles de requete/reponse (Pydantic)
|   |   +-- routers/           # Un fichier par groupe d'endpoints
|   |   |   +-- health.py        validation.py   prediction.py
|   |   |   +-- explain.py       scenario.py     monte_carlo.py
|   |   |   +-- metadata.py      chat.py
|   |   +-- services/          # Logique metier decouplee du HTTP
|   |       +-- artifacts.py          # Chargement/cache de l'artefact joblib
|   |       +-- model_service.py      # Feature engineering + predict_proba
|   |       +-- shap_service.py       # TreeExplainer SHAP
|   |       +-- scenario_service.py   # PD baseline vs. modifiee (delta_pd)
|   |       +-- monte_carlo_service.py# Stress test stochastique vectorise
|   |       +-- llm_service.py        # Appel async httpx vers Ollama
|   |       +-- feature_engineering.py# Ratios derives a l'inference
|   |       +-- validation_service.py # Controle des features
|   |       +-- eda_service.py        # Resume exploratoire du dataset
|   +-- requirements.txt
|   +-- .env / .env.example
+-- frontend/                  # Tableau de bord statique (HTML/CSS/JS)
|   +-- index.html  app.js  styles.css
+-- data/                      # Datasets + artefacts modeles (.joblib, .json)
+-- scripts/                   # Entrainement, generation de donnees, IFRS9
|   +-- train_option_c_xgboost.py
|   +-- train_option_c_randomforest.py
+-- notebook/                  # Notebooks de modelisation (Options A/B/C)
+-- dataset_IFRS9/             # Tables IFRS9 (PD/LGD/EAD/discount/exposure)
+-- notebook_ifrs/             # Notebooks et synthese ECL
+-- texts/                     # Rapports, dictionnaire de variables
\end{lstlisting}

Ce découpage applique le principe de \textbf{séparation entre la couche HTTP et la
logique métier} : les fichiers de \texttt{routers/} ne contiennent que la déclaration
des points d'entrée et la validation des entrées/sorties, tandis que toute la logique
analytique réside dans \texttt{services/}. Un service peut ainsi être réutilisé par
plusieurs endpoints (par exemple \texttt{model\_service.predict} est appelé à la fois
par \texttt{/predict}, \texttt{/scenario} et \texttt{/chat/analyze}).

% ────────────────────────────────────────────────────────────
\section{Choix technologiques et environnement de développement}

Le tableau~\ref{tab:techno} récapitule la pile technologique retenue et la
justification de chaque choix au regard des besoins du projet.

\begin{table}[h!]
\centering
\caption{Pile technologique de la plateforme}
\label{tab:techno}
\renewcommand{\arraystretch}{1.3}
\begin{tabularx}{\textwidth}{@{}l l Y@{}}
\toprule
\textbf{Domaine} & \textbf{Technologie} & \textbf{Justification} \\
\midrule
Langage         & Python 3            & Écosystème data science / ML mature \\
API REST        & FastAPI + Uvicorn   & Asynchrone, typage natif, OpenAPI/Swagger auto \\
Validation      & Pydantic v2         & Schémas typés, configuration par variables d'env. \\
ML --- arbres   & XGBoost, scikit-learn (RandomForest) & Modèles ensemblistes performants \\
Prétraitement   & scikit-learn        & \texttt{ColumnTransformer}, encodeurs, scalers \\
Explicabilité   & SHAP (\texttt{TreeExplainer}) & Explications cohérentes globales et locales \\
Sérialisation   & joblib              & Persistance préprocesseur + modèle \\
LLM             & Ollama (API compatible OpenAI) & Inférence locale, souveraineté des données \\
Client HTTP     & httpx (async)       & Appels non bloquants vers le LLM \\
Frontend        & HTML/CSS/JS, Chart.js & Tableau de bord sans étape de build \\
\bottomrule
\end{tabularx}
\end{table}

\subsection{Configuration par variables d'environnement}

Toute la configuration (chemins des artefacts, seuil de décision, activation et
paramétrage du LLM, hôte/port) est centralisée dans \texttt{config.py} via
\texttt{BaseSettings} de Pydantic, alimentée par un fichier \texttt{.env}. Ce
mécanisme découple le code de l'infrastructure : il permet de basculer de modèle
(RandomForest $\leftrightarrow$ XGBoost), d'activer ou non le LLM, ou de durcir le
seuil de décision sans modifier le code.

\begin{lstlisting}[caption={Extrait de configuration (.env)}, label={lst:env}]
MODEL_ARTIFACT_PATH=d:/BFPME/data/option_c_randomforest_model.joblib
MODEL_METRICS_PATH=d:/BFPME/data/option_c_randomforest_metrics.json
DEFAULT_THRESHOLD=0.5            # PD >= seuil => defaut predit
LLM_ENABLED=true
LLM_BASE_URL=http://127.0.0.1:11434/v1   # Ollama
LLM_MODEL=qwen2.5:7b
LLM_MAX_TOKENS=500
\end{lstlisting}

% ────────────────────────────────────────────────────────────
\section{Conception des données}

\subsection{Sources et nature des données bancaires}

Les données proviennent des dossiers de crédit PME : informations signalétiques de
l'entreprise (ancienneté, secteur, région, structure juridique, classification PME,
effectif), états financiers (chiffre d'affaires, total actif/passif, capitaux
propres, dettes court et long terme, fonds de roulement, BFR, trésorerie), variables
de crédit (montant du prêt, durée, type de crédit, objectif, type de taux, valeur des
garanties), variables comportementales (retards de paiement, impayés) et variables
de scoring interne BFPME. La variable cible est \texttt{default\_flag} (défaut /
non-défaut).

Conformément au cadre du projet, le jeu de données exploité
(\texttt{final\_dataset.csv}) est \textbf{synthétique} mais métier-cohérent : il est
généré à partir d'un \textit{driver} latent de risque dérivé du score BFPME, de la
classe de risque et du type de projet (création / extension), de sorte que les
variables financières, comportementales et de garantie restent corrélées de façon
réaliste. Il sert de support de prototypage, de comparaison de modèles et
d'ingénierie de variables.

\subsection{Modèle de données et entités principales}

Le modèle de données structure les entités centrales du domaine du crédit PME et
leurs relations. La figure~\ref{fig:classes_donnees} (diagramme de classes, non
modifié) en formalise la structure ; les entités principales sont :

\begin{itemize}
  \item \textbf{Entreprise / PME} : caractéristiques signalétiques et financières du
        demandeur ;
  \item \textbf{Dossier de crédit} : montant, durée, type, objectif, garanties ;
  \item \textbf{Décision} : PD estimée, classe de risque, seuil appliqué, explication
        associée ;
  \item \textbf{Exposition IFRS~9} : paramètres PD / LGD / EAD et résultats ECL
        (module réglementaire).
\end{itemize}

\subsection{Pipeline de prétraitement (artefact joblib)}

Le préprocesseur et le modèle sont entraînés conjointement puis sérialisés dans un
\textbf{unique artefact joblib} contenant : le \texttt{preprocessor}
(\texttt{ColumnTransformer} ajusté), le \texttt{model} entraîné et la liste
\texttt{allowed\_columns} des colonnes attendues. Le préprocesseur applique, selon le
type de variable, la chaîne suivante :

\begin{enumerate}
  \item \textbf{Variables numériques sensibles aux valeurs extrêmes} : imputation par
        la médiane $\rightarrow$ \textbf{Winsorisation} (écrêtage aux centiles 1 et 99)
        $\rightarrow$ \textbf{standardisation} ;
  \item \textbf{Variables numériques simples} : imputation médiane $\rightarrow$
        standardisation ;
  \item \textbf{Variables ordinales} (ex. durée de crédit) : imputation par mode
        $\rightarrow$ \texttt{OrdinalEncoder} avec ordre métier explicite ;
  \item \textbf{Variables nominales} (secteur, région, structure juridique, type de
        crédit, scénario économique, \dots) : imputation par mode $\rightarrow$
        \texttt{OneHotEncoder} (\texttt{handle\_unknown=ignore}).
\end{enumerate}

Ce découpage par type de variable, encapsulé dans un \texttt{ColumnTransformer},
porte la dimensionnalité de \textbf{52~variables brutes à 80~variables transformées}.
La figure~\ref{fig:pipeline_prepro} en illustre le déroulé.

\begin{figure}[h!]
  \centering
  \includegraphics[width=0.95\textwidth]{pipeline_preprocessing}
  \caption{Pipeline de prétraitement sérialisé dans l'artefact joblib}
  \label{fig:pipeline_prepro}
\end{figure}

\subsection{Ingénierie de variables à l'inférence}

Un choix de conception déterminant consiste à appliquer l'\textbf{ingénierie de
variables avant le préprocesseur, à l'inférence}, et non à la figer dans l'artefact.
Le service \texttt{feature\_engineering} dérive à la volée des ratios métier
discriminants à partir des champs bruts du dossier, par exemple :

\[
\text{loan\_to\_revenue} = \frac{\text{montant\_pret}}{\text{chiffre\_affaires\_annuel}},
\qquad
\text{cash\_to\_debt} = \frac{\text{tresorerie\_disponible}}{\text{dette\_totale}}
\]
\[
\text{payment\_stress\_index} =
\mathrm{clip}\!\left(\tfrac{\text{retards\_jours}}{90},0,3\right) +
\mathrm{clip}\!\left(\text{impayes\_to\_loan},0,3\right),
\qquad
\text{garantie\_to\_loan} = \frac{\text{valeur\_garanties}}{\text{montant\_pret}}
\]

D'autres ratios sont calculés selon la même logique (\texttt{debt\_total},
\texttt{impayes\_to\_loan}, \texttt{fdr\_bfr\_pressure},
\texttt{financement\_externe\_intensity}). Cette approche garantit que la
\textbf{même logique de dérivation} est appliquée à l'entraînement, à la prédiction
unitaire, au scénario et à la simulation Monte~Carlo, évitant toute dérive entre
apprentissage et production.

\subsection{Données du module IFRS 9}

Le module IFRS~9, découplé de l'API, s'appuie sur un ensemble de tables structurées
(\texttt{dataset\_IFRS9/}) : instantané d'exposition, scénarios macro-économiques,
courbes PD, LGD, EAD et facteur d'actualisation. Le calcul de l'\textit{Expected
Credit Loss} y est conduit selon :
\[
  \mathrm{ECL}_t = \mathrm{PD}_t \times \mathrm{LGD}_t \times \mathrm{EAD}_t \times \mathrm{DF}_t,
  \qquad
  \mathrm{ECL}_{\text{pondérée}} = \sum_{s} w_s \times \mathrm{ECL}^{\text{lifetime}}_s
\]
où la PD issue de la modélisation prédictive alimente directement le moteur de calcul,
et où la pondération forward-looking par scénario assure la conformité réglementaire.

% ────────────────────────────────────────────────────────────
\section{Conception de l'API et des services}

\subsection{Points d'entrée REST}

L'API expose un ensemble cohérent de points d'entrée, chacun associé à un routeur
dédié. Le tableau~\ref{tab:endpoints} en présente la cartographie.

\begin{table}[h!]
\centering
\caption{Cartographie des points d'entrée de l'API}
\label{tab:endpoints}
\renewcommand{\arraystretch}{1.3}
\begin{tabularx}{\textwidth}{@{}l l Y@{}}
\toprule
\textbf{Méthode} & \textbf{Endpoint} & \textbf{Rôle} \\
\midrule
GET  & \texttt{/health}              & Sonde de disponibilité du service \\
GET  & \texttt{/metadata/model}      & Métadonnées et métriques du modèle \\
GET  & \texttt{/metadata/eda}        & Résumé exploratoire du dataset (graphes) \\
POST & \texttt{/metadata/reload-model} & Rechargement de l'artefact sans redémarrage \\
POST & \texttt{/validate}            & Contrôle de cohérence du dossier \\
POST & \texttt{/predict}             & Prédiction de la PD et de la classe de risque \\
POST & \texttt{/explain}             & Contributions SHAP (top-$k$) \\
POST & \texttt{/scenario}            & Simulation what-if (PD baseline vs. modifiée) \\
POST & \texttt{/monte-carlo}         & Test de stress stochastique \\
POST & \texttt{/chat/analyze}        & Orchestration complète + synthèse LLM \\
\bottomrule
\end{tabularx}
\end{table}

\subsection{Cycle de vie de l'artefact}

L'artefact est chargé \textbf{une seule fois} au démarrage de l'application
(événement \texttt{startup}) et mis en cache en mémoire (\texttt{lru\_cache}). Le
point d'entrée \texttt{/metadata/reload-model} permet de vider ce cache et de
recharger le modèle à chaud, sans interruption de service, après un réentraînement.
Cette conception répond à l'exigence de performance (pas de rechargement disque à
chaque requête) tout en préservant la maintenabilité opérationnelle.

% ────────────────────────────────────────────────────────────
\section{Architecture de l'agent IA et orchestration}

\subsection{Pipeline d'orchestration}

Le c\oe ur de l'agent réside dans l'endpoint \texttt{/chat/analyze}, qui orchestre
les outils analytiques en fonction des indicateurs portés par la requête
(\texttt{run\_validation}, \texttt{run\_prediction}, \texttt{run\_explanation},
\texttt{scenario\_changes}). Il compose les résultats dans une structure unique
\texttt{tool\_results}, puis délègue la rédaction de la synthèse au LLM. La
figure~\ref{fig:seq_chat} (diagramme de séquence, non modifié) en détaille les
échanges.

\begin{figure}[h!]
  \centering
  \includegraphics[width=0.95\textwidth]{sequence_chat_analyze}
  \caption{Diagramme de séquence de l'orchestration \texttt{/chat/analyze}}
  \label{fig:seq_chat}
\end{figure}

\begin{lstlisting}[caption={Logique d'orchestration (extrait)}, label={lst:orch}]
if run_validation:  tool_results["validation"]  = validate_features(features)
if run_prediction:  tool_results["prediction"]  = predict(features, threshold)
if run_explanation: tool_results["explanation"] = explain_prediction(features, k)
if scenario_changes: tool_results["scenario"]   = run_scenario(features, changes)
# synthese narrative deleguee au LLM (si active), sinon sorties brutes
\end{lstlisting}

\subsection{Seuils de confiance et classes de risque}

L'agent applique un \textbf{seuil de décision} configurable (\texttt{DEFAULT\_THRESHOLD},
0{,}5 par défaut) : la classe prédite vaut « défaut » dès que la PD dépasse ce seuil.
Indépendamment, une classe de risque qualitative est attribuée à des fins d'escalade :
\textit{low} (PD~$<$~0{,}4), \textit{medium} ($0{,}4 \le$~PD~$<$~0{,}7) et
\textit{high} (PD~$\ge$~0{,}7). Ce double mécanisme — seuil quantitatif et classe
qualitative — matérialise la logique \textit{human-in-the-loop} : les dossiers à
risque élevé sont signalés pour escalade vers le comité de crédit, l'IA assistant la
décision sans s'y substituer.

\subsection{Séparation stricte des rôles modèle / LLM}

Un principe de conception fondamental gouverne l'intégration du LLM : la
\textbf{séparation stricte des responsabilités}. Le modèle prédictif estime le risque ;
le LLM se limite à \textbf{expliquer et reformuler} en langage naturel à partir des
sorties calculées. Le \textit{prompt} système impose explicitement de n'utiliser que
les sorties fournies et de ne jamais inventer, arrondir ou estimer une valeur
numérique. L'appel est asynchrone (\texttt{httpx}), à température basse (0{,}1) pour
limiter la variabilité, et toute défaillance du LLM est \textbf{interceptée} sans
jamais invalider les résultats déjà calculés — l'analyste reçoit alors les sorties
brutes du modèle, garantissant la continuité de service.

% ────────────────────────────────────────────────────────────
\section{Module d'explicabilité SHAP}

Le service d'explicabilité instancie un \texttt{TreeExplainer} (mis en cache) sur le
modèle entraîné et calcule les valeurs SHAP sur la \textbf{matrice de variables
transformées}. Pour chaque prédiction, il extrait les $k$ contributions de plus forte
amplitude, indique pour chacune sa direction (\textit{increase\_risk} /
\textit{decrease\_risk}) et restitue la valeur de base (\textit{base value}). Cette
restitution sert à la fois l'explicabilité \textbf{locale} (justification d'une
décision individuelle) et, par agrégation, l'explicabilité \textbf{globale}
(importance des variables sur le portefeuille), répondant aux besoins respectifs des
analystes et des gestionnaires de risque.

% ────────────────────────────────────────────────────────────
\section{Moteur de simulation Monte Carlo}

Le moteur de stress test génère, à partir d'un client de base et d'un ensemble de
\textbf{chocs} sur des variables choisies, $N$ scénarios perturbés, puis exécute le
modèle sur l'ensemble du lot en \textbf{une seule passe vectorisée}
(un \texttt{transform} + un \texttt{predict\_proba}), ce qui garantit la performance.
Les chocs sont appliqués aux variables \emph{brutes} (selon une loi normale ou
uniforme), puis les ratios d'ingénierie sont re-dérivés pour rester cohérents.

La distribution de PD obtenue est synthétisée par des \textbf{indicateurs de risque}
directement exploitables : moyenne et écart-type, percentiles, \textbf{Value-at-Risk
à 95~\%} (\textit{VaR}$_{95}$), \textbf{Expected Shortfall} (perte moyenne au-delà de
la VaR), probabilité de dépassement du seuil, répartition par classe de risque,
histogramme et scénarios extrêmes. Ces mesures établissent le lien opérationnel avec
la dimension forward-looking de l'IFRS~9 décrite au chapitre précédent.

% ────────────────────────────────────────────────────────────
\section{Modélisation UML}

La conception du système a été formalisée à l'aide du langage UML. Les diagrammes
présentés ci-dessous \textbf{ne sont pas modifiés} et sont repris tels qu'élaborés
lors de la phase de conception.

\subsection{Diagramme de cas d'utilisation}

Le diagramme de cas d'utilisation (figure~\ref{fig:uc}) identifie les acteurs
principaux — l'analyste crédit, le gestionnaire de risque et le comité de crédit —
et leurs interactions avec la plateforme : prédire la PD, expliquer une décision,
simuler un scénario, lancer un test de stress et dialoguer en langage naturel.

\begin{figure}[h!]
  \centering
  \includegraphics[width=0.9\textwidth]{diagramme_cas_utilisation}
  \caption{Diagramme de cas d'utilisation de la plateforme}
  \label{fig:uc}
\end{figure}

\subsection{Diagramme de classes}

Le diagramme de classes (figure~\ref{fig:classes_donnees}) formalise les entités du
domaine et les composants de service (préprocesseur, modèle, explicateur,
orchestrateur) ainsi que leurs relations.

\begin{figure}[h!]
  \centering
  \includegraphics[width=0.95\textwidth]{diagramme_classes}
  \caption{Diagramme de classes du système}
  \label{fig:classes_donnees}
\end{figure}

\subsection{Diagramme de déploiement}

Le diagramme de déploiement (figure~\ref{fig:deploiement}) situe les n\oe uds
d'exécution : le serveur applicatif (FastAPI/Uvicorn), le serveur d'inférence LLM
local (Ollama) et le navigateur client hébergeant le tableau de bord.

\begin{figure}[h!]
  \centering
  \includegraphics[width=0.9\textwidth]{diagramme_deploiement}
  \caption{Diagramme de déploiement de la plateforme}
  \label{fig:deploiement}
\end{figure}

% ────────────────────────────────────────────────────────────
\section{Planification des sprints et des releases}

Conformément au choix méthodologique SCRUM (chapitre~1), le développement a été
organisé en releases incrémentales, chacune livrant un sous-ensemble cohérent et
validable de la plateforme. Le tableau~\ref{tab:releases} synthétise cette
planification.

\begin{table}[h!]
\centering
\caption{Planification des releases du projet}
\label{tab:releases}
\renewcommand{\arraystretch}{1.3}
\begin{tabularx}{\textwidth}{@{}l Y l@{}}
\toprule
\textbf{Release} & \textbf{Périmètre} & \textbf{Chapitre} \\
\midrule
Release~I   & Compréhension métier et ingénierie des données (cycle du crédit, collecte, nettoyage, \textit{feature engineering}) & Ch.~4 \\
Release~II  & Modélisation prédictive (régression logistique, Random Forest, Gradient Boosting) et explicabilité SHAP & Ch.~5 \\
Release~III & Agent IA, orchestration décisionnelle et intégration LLM en langage naturel & Ch.~6 \\
Release~IV  & Tests de stress Monte~Carlo et conformité IFRS~9 (ECL, staging, forward-looking) & Ch.~7 \\
\bottomrule
\end{tabularx}
\end{table}

Chaque release a été clôturée par une revue de sprint avec les parties prenantes de
MPSoft, suivie d'une rétrospective, garantissant un rythme de livraison régulier et
une adaptation continue aux retours métier.

% ────────────────────────────────────────────────────────────
\section*{Conclusion}

Ce chapitre a présenté la conception et l'architecture technique de la plateforme :
une architecture en cinq couches matérialisée par une structure de projet claire
séparant la couche HTTP de la logique métier ; une pile technologique (FastAPI,
scikit-learn/XGBoost, SHAP, Ollama) cohérente avec les exigences de performance,
d'explicabilité et de souveraineté des données ; une conception des données
articulant un préprocesseur sérialisé, une ingénierie de variables à l'inférence et
les tables IFRS~9 ; et une orchestration d'agent IA assurant une séparation stricte
entre estimation du risque et restitution en langage naturel. La modélisation UML et
la planification des sprints complètent cette vue de conception. Les chapitres
suivants détaillent la réalisation de chacune des releases, en commençant par la
compréhension métier et l'ingénierie des données.

\end{document}
